% !TeX program = xelatex
\documentclass{UAmathtalk}
\usepackage{setspace}
\renewcommand{\when}{12:00 p.m.}
\renewcommand{\where}{Hudson 0110}
%\renewcommand{\where}{Hudson 0110.\\ \href{https://albany.zoom.us/j/81219478289?pwd=N5fy7vxKBXqlQ4mAq4KaeaAM5iAHi4.1}{Zoom ID: 812\,1947\,8289 - Passcode: 233444}}
\author{David Hien}
\address{Rutgers University}
\urladdr{https://davidhien.github.io/}
\title{Cycling Signatures: Analyzing Recurrence in Dynamical Systems via~Topological~Data~Analysis}
\date{Friday, October 30th, 2026}


\begin{document}

\maketitle
%\renewcommand*{\abstractsize}{\normalsize}
\renewcommand*{\abstractsize}{\fontsize{13.5}{16.2}\selectfont}
\begin{abstract}
    {
Recurrence is a fundamental characteristic of dynamical systems with complicated behavior. While simple forms of recurrence, such as periodic orbits, admit clear qualitative descriptions, more complex recurrent dynamics are difficult to characterize, particularly in high-dimensional systems and in the presence of noise.

Cycling signatures provide a coarse description of recurrent dynamics in terms of elementary building blocks called cycling motions. They are computable from from time series data and associate a filtered vector space to each trajectory segment which describes the cycling motions represented in that segment at different spatial scales. 
Statistics of these signatures reveal how cycling motions occur and how the dynamics moves between them.

In this talk, we introduce cycling signatures, address some computational aspects, and illustrate the information they capture. We then discuss distributions of cycling signatures and use Birkhoff’s ergodic theorem to relate time averages along trajectories to statistics of the underlying invariant measure.     }
\end{abstract}
\end{document}


